\chapter{Introduction}
\minitoc

\section{Classical and Quantum Computing}
% What is QM
Quantum mechanics (QM) successfully describes the behaviour of sufficiently small things.
The theory is indispensable in the description of sub-atomic particles, atoms, the
molecule they compose, crystals, light, the vacuum, and many other parts of our
reality. That is to say, given a choice, bet things obey quantum mechanics. 
It was recognised in the 1980s that quantum mechanics and information must be unified for fully unlocking the burgeoning field of computation~\cite{benioff_computer_1980, feynman_simulating_1982}. Motivated with the intuition that if reality is described by QM, then surely
quantum objects will offer more efficient strategies to simulate reality, the idea of the \emph{Quantum Comptuter} was formed. We will elaborate this concept, but first 
we donate a paragraph to the classical computers that have gone before.\\
% Classical Computing
What is the goal of computing? We will define ``compute'' as modelling some physical system with another,
disparate, physical system. This could be an analogue simulator with continuous
input and outputs, or it could be some digital model that mimics the system at
an arbitrary precision. Apart from the engineering practicalities of realising
the device, this analogue and digital distinction may be abstracted away. The
\emph{Computer} enters when Turing and Church proposed a class of these simulators that not
only mimic one specific system, but rather can model any physical system,
provided enough time and memory~\cite{turing_computable_1937, church_unsolvable_1936}. A universal simulator is of course a
useful tool. With sufficient understanding of the physical system we wish to model, we
may observe \emph{simulated} behaviour without requiring control over the actual
system. 
A profound result in the study of digital universal simulators is the small set of elementary operations required for this universality, for example, the boolean set \{AND,OR,NOT\}, are sufficient to simulate any function to an arbitrary precision.\\
% Why we need QC
However, \emph{universal} does not mean \emph{practical} ---  although a physical device
is able to simulate anything, it may still not be feasible for an operator to ``code''
some desired model\footnote{play with esolangs for a few hours to confirm this
for yourself.}, or for a useful result to be returned in reasonable
time --- i.e. is tractable\footnote{factor a 2048-bit number for a few billion years to confirm this for
yourself.}. Regardless, highly advanced and operable computers have been
constructed and have successfully modelled and transformed almost all aspects
of modern life. \\
Unfortunately, many problems intractable using classical computers are of great
practical importance, for example accurately modelling chemical systems for
industrial, and drug design, and simulating solid-state systems for fundamental
material, and physical research.\\
Both these examples intentionally fall into the simulation of quantum systems.
Models of quantum systems are notoriously expensive to encode, 
fundamentally due to
the concepts of superposition and interference --- a pure quantum state is generally
described by the superposition of all possible $d$ basis states, where $d$ is
the system dimension. Given an underlying number, $n$, of $q$-level subsystems that compose the system, we find the dimension exponentially scales as $d=q^n$. Each basis state of the $n$-body system will have some
associated amplitude and phase, resulting in $2q^n-2$ real
parameters required to track and accurately model the pure quantum state. This
exponential scaling is not kind to a digital computer.\\
% What is QC
Returning to the 1980s, it was realised that this exponential scaling could provide the very resources required to simulate other quantum systems. Can we construct a Turing
machine that utilises the resource scaling of a quantum system~\cite{deutsch_quantum_1985}\footnote{As an
aside, the implications of quantum mechanics on information
extend further than purely mimicking other dynamic systems -- using and
distributing quantum objects also allows for dominant strategies in
cryptography~\cite{bennett_quantum_2014}, and some competitive (although currently esoteric)
games~\cite{drmota_experimental_2025}.} --- a quantum computer (QC)? \\
The ingredients for such a device are well
known~\cite{divincenzo_physical_2000}, and have been demonstrated on multiple
physical platforms to date. Theoretical work is placing ever-more lenient
constraints on resource requirements, as algorithmic strategies are developed,
while experimental work is pushing to scale the avaiable resources, and
performance of actual hardware. The two frontiers are yet to meet at a
realisation of a useful universal QC, however the list of prerequirements is
shortening rapidly.
 We shall give a short overview of the distinct platforms in
development for scaling QC, and then concentrate on trapped-ions  --- the
experimental platfrom we use in this thesis.

\section{Quantum Computing Platforms}

An effective quantum computer requires meticulous control over a collection of
simple quantum systems. A good starting point is thus selecting that simple
system to build our control around. Nature provides us with a bounty of choices
in this regard: atoms. It is definitively difficult to choose a simpler and more reproducible object: atoms of the same species are
identical to one another no matter where or when you are.
We may use their electronic energy structure, which atomic physicists have
scrutinized for over a century, to define our quantum bits, and we can use an
assortment of electromagnetic frequencies to manipulate these internal states.
This is the strategy of both the ``neutral-atom'', and ``trapped-ion''
platforms, with the latter choosing to remove or add a few electrons to more
easily manipulate the position of the atoms in space\footnote{We note that the
neutral-atom platform has made significant progress in the spatial control of atoms in optical tweezer traps, however, the trap depths between the two platforms is incomparable.}. The trapped-ion platform also has access to the motion of the ions, which may be used as continuous variable quantum resources --- so called qumodes --- that
can be coherently controlled, and initialised to highly non-classical states.\\
Photons are another ``natural'' choice for the quantum-resource. They may either be treated also as continuous quantum systems, or have some discrete level system encoded, such as a qubit. The photonic qubit is unique in being a ``flying-qubit'': photons are
able to transmit information over long distances, lending them particular interest in their use for interconnecting distant quantum computers.\\
The final class of platforms is that of
``manufactured'' quantum systems, such as superconducting circuits, quantum dots, and NV centers;
benefiting from control over the level structure, and owing to their solid
state, presumably allowing for easier integration into existing semiconductor
fabrication techniques.\\
At this time, it is unclear whether one of these platforms,
or an entirely new architecture, will be the first to reach a scaled device.
Each have their own advantages and unique technical challenges, and each are in
rapid active development. The next section outlines some technical challenges in
scaling current quantum hardware, focussing on the trapped-ion platform studied in this thesis.

\section{Scaling Trapped Ion Systems}

For useful quantum computation and simulation, we require access to a sufficiently large Hilbert space,
with sufficiently accurate control with which to traverse it. How do we scale the ion-trap
platform to provide those resources? In principle, we know \emph{how}, but the required engineering and control presents unique and technical challenges. I identify three
specific challenges associated to scaling, that motivate the work presented in this thesis:
increasing the available quantum resources, controlling those resources, and characterising and mitigating the errors that arise as the system grows.\\

\subsubsection{Challenge I: Scaling Quantum Resources}
We need larger Hilbert spaces --- to solve a ``useful'' computation, such as
Shor's algorithm~\cite{shor_polynomial-time_1997} applied to factoring RSA-2048
integers, it is currently estimated that
$10^4-10^5$ physical qubits will be required~\cite{cain_shors_2026,
webster_pinnacle_2026}. Architectures for such scaling include quantum charge-coupled device (QCCD)~\cite{kielpinski_architecture_2002},
to reconfigure an array of qubits using multi-zone surface traps; and optical
networked modules~\cite{luo_protocols_2009}, using photonic interconnects to combine many small
qubit registers into a single traversable Hilbert space. Another disparate strategy
--- continuous variable quantum computing (CVQC) --- replaces many low-dimensional
qubits, with few high-dimensional qumodes. Qumodes store information in the continuous states of a quantum harmonic oscillator, as such, in principle have an infinite-dimensional Hilbert space. This exchanges the
technical issues of controlling vast numbers of qubits, to instead requiring
truly exquisite control of a few qumodes. The innate utility of this continuous
system, is exemplified by a CVQC scheme for
factoring RSA-2048 integers using only three qumodes and a single
qubit~\cite{brenner_factoring_2024}. Although this particular protocol is
unfeasible due to technically prohibitive required high-energy
oscillator states, scaling accessible Hilbert spaces via qumodes remains a promising
strategy. Trapped ions therefore provide two complementary routes to increasing the accessible Hilbert space: increasing the total number of ions in the register, and utilising the motional modes of the ions already present within that register.

\subsubsection{Challenge II: Scaling Control}
Access to a sufficiently large Hilbert
space is useful only if we can reliably and selectively navigate it.
What are the elementary operations we require, and how do we
physically implement them on our system? The elementary unitary operations often
are expressed as minimal gate-sets allowing for universal quantum computing.
These are typically comprised of one, and two-body
interactions~\cite{divincenzo_two-bit_1995,
kitaev_quantum_1997,braunstein_quantum_2005,lloyd_quantum_1999}, which greatly
simplifies the elementary operation, and the required number of implementable
interactions. This does comes at the considerable cost of
``addressability'' --- we must be able to selectively drive interactions locally
on either one, or two physical carriers. For trapped-ion qubits, this requires a
form of spatial selectivity, while for motional-state qumodes, often requires a combination of
frequency and spatial selectivity. This technical constraint is that of
``control-signal delivery'', a common problem to both classical and quantum
computing. Trapped-ion systems require hardware for delivering spatially, and spectrally selective electronic control, and optical-beam
delivery, both of which are under active development to enable
scalability~\cite{malinowski_how_2023,mehta_integrated_2020}.\\

\subsubsection{Challenge III: Scaling with Imperfect Systems}
The last specific challenge we discuss is that of unwanted evolution due to \emph{noise}.
Qubits stored in the electronic states of ions are sensitive to external electric
and magnetic fields, while qumodes are particularly sensitive to electric fields originating from the trap electrodes. The effects of noise on the desired
state evolution are varied, but
if left unchecked, inevitably lead to information loss. 
Furthermore, increased system scale and complexity changes the relevant error
processes. In particular, noise leading to unintended correlations across the register, and across time can be difficult to quantify and mitigate.
To combat this noise, we first require methods to characterise performance, and
quantify how this noise leads to logical errors in an executed algorithm.
This motivates the wide field of
quantum characterisation and verification~\cite{eisert_quantum_2020,
hashim_benchmarking_2023}. A varied collection of characterisation protocols
exist due to an unavoidable trade-off between \emph{expressibility} and
\emph{scalability} --- The systems we wish to characterise are inherently large,
and thus full descriptions require vast amounts of information. So we often sit
in two regimes: we may either fully characterise few-body
interactions~\cite{chuang_prescription_1997}, or we must use abstracted, and
system-wide metrics to gauge performance of the scaled
system~\cite{erhard_characterizing_2019,mckay_benchmarking_2023}.\\
When noise is suitably suppressed
through hardware engineering, we may then utilise algorithmic techniques to be
robust to the remaining imperfections. Foremost, analogous to classical
computing, it is possible to detect and correct errors during the execution of
an algorithm~\cite{aharonov_fault-tolerant_2006}. This comes at a considerable
overhead in both register size and circuit depth, only recently in reach of
experimental platforms, but only for small numbers of logical qubits~\cite{steane_overhead_2003}. A
complimentary method for suppressing the impact of errors is through myriad
error mitigation techniques, using repeat measurements to infer error-suppressed
algorithmic outputs~\cite{cai_quantum_2023}. \\
Therefore, even with imperfect physical qubits or qumodes, it appears to still be physically possible to construct useful, universal quantum computers.


\section{Structure of this Thesis}

    This thesis describes the development of a trapped-ion platform for
    multi-ion spin–motion control, and its application to the characterisation
    of correlated errors, hybrid spin–oscillator control, and quantum error
    mitigation.\\

    \noindent Chapter~\ref{ch:theory} provides further details on the trapped-ion platform, and the underlying light-matter interactions enabling the use of ions as quantum-information carriers. \\
    Chapter~\ref{ch:apparatus} details the ion-trap 
    apparatus constructed for this thesis, enabling the study of multi-qubit and
    qumode registers.\\
    Chapter~\ref{ch:characterisation} presents the hardware characterisation, giving details on experimental calibrations, the initialisation of the qubits and qumodes, and elementary single- and two-qubit gate performance. \\
    Chapter~\ref{ch:k-copy} presents a novel and scalable randomised benchmarking method
    for characterising and quantifying spatially correlated errors in multi-qubit
    devices, including experimental results on up-to an 8-ion chain.\\
    Chapter~\ref{ch:QEM} explores two ongoing research directions enabled by the
    newly constructed apparatus. First, we outline how to utilise the qumodes in our system, providing a gate-set for full qubit-qumode control. We discuss a practical way to drive this full gate-set, and by this method, demonstrate spin-conditioned motional squeezing. Second, we experimentally investigate virtual distillation, an error mitigation technique with
    practical application to computation involving both spin and motional
    degrees of freedom.\\ 
    Chapter~\ref{ch:conclusion} concludes this thesis by summarising its contributions, and providing an outlook for future research enabled by the experimental platform.\\